\section{Positive capacity implies reconstruction}

This direction uses only Lemma~\ref{lem:local}. We introduce a harmless
optional erasure to ensure nonnegative information deficits without
assuming Proposition~\ref{prop:input}.

\subsection{Products followed by erasure}

For symmetric experiments with information $u,v$, form their raw product,
whose information is $m_*$. If $m_*>u+v$, retain the combined observation
with probability $(u+v)/m_*$ and otherwise output a fixed erasure symbol;
if $m_*\le u+v$, retain it always. The retention coin is independent
of the spin and of all observations; its parameter depends only on the
experiment laws. Conditional on erasure, the posterior is uniform.
Thus the new experiment has information
\[
m'=\min\{m_*,u+v\}
\]
and third-moment magnitude at most that of the raw product. Lemma
\ref{lem:local} gives
\begin{equation}\label{eq:erasure}
0\le u+v-m'\le5(h_1v+u h_2),\qquad
h'\le h_1+h_2+376uv.
\end{equation}
The experiment remains symmetric and is a degradation of the original
pair. Independent random coins in disjoint subtrees preserve the
conditional independence needed at subsequent products.

\subsection{A flow profile and finite observations}

Take an admissible flow of mass $\alpha>0$, and scale it down so that
$\alpha\le1$. For $0<\varepsilon\le1$, set
\[
\Theta(v)=\varepsilon\theta(e_v)\ (v\ne o),\qquad
\Theta(o)=M=\varepsilon\alpha,\qquad
X_v=\frac{\Theta(v)}{\rho(v)}.
\]
Flow conservation says
\begin{equation}\label{eq:profile}
X_v=\sum_{w\in\ch(v)}\lambda_{vw}^2X_w.
\end{equation}
Every positive-flow edge extends to an infinite ray, so admissibility
gives $X_v\le\varepsilon$ for $v\ne o$. Also $X_o=M\le\varepsilon$.
On every ray,
\begin{equation}\label{eq:energy}
\sum_{v\in\xi}X_v^2\le2\varepsilon^2.
\end{equation}
This normalization replaces the uniform degree bound used to control
the root in the homogeneous argument.

Fix a depth $n$. At each $v\in L_n$, pass its spin through an additional
independent channel of parameter $t_v=\sqrt{X_v/3}$. The terminal
experiment then has
\begin{equation}\label{eq:terminal}
\widehat m_v=X_v,\qquad
\widehat h_v=(X_v/3)^{3/2}\le X_v/3.
\end{equation}
Proceed upward. At each vertex pass each constructed child experiment
through its actual edge channel, and combine the branches sequentially
using \eqref{eq:erasure}. All experiments so obtained are degradations
of the level-$n$ observation. Their information satisfies
\begin{equation}\label{eq:below}
0\le\widehat m_v\le X_v
\end{equation}
by induction and \eqref{eq:profile}. Empty observations and zero-flow
subtrees cause no difficulty.

\subsection{Local deficits with constants independent of degree}

At an internal vertex, write
\[
b_w=\lambda_{vw}^2\widehat m_w,\qquad
h'_w=\lambda_{vw}^3\widehat h_w,\qquad
S=\sum_w b_w\le X_v,\qquad H=\sum_w h'_w.
\]
Let
\begin{equation}\label{eq:deficit}
D_v=\sum_w\lambda_{vw}^2\widehat m_w-\widehat m_v.
\end{equation}
With $C_0=376$, sequential use of \eqref{eq:erasure} gives
\begin{equation}\label{eq:sharpaggregate}
\widehat h_v\le H+\frac{C_0}{2}S^2,
\qquad 0\le D_v\le5SH+\frac{5C_0}{6}S^3.
\end{equation}
Here are the details of the degree-independent sums. Order the children
and put $s_{j-1}=\sum_{i<j}b_i$, $H_{j-1}=\sum_{i<j}h'_i$.
The partial product has $m\le s_{j-1}$ and
$h\le H_{j-1}+(C_0/2)s_{j-1}^2$. The accumulated third-moment
errors are bounded by
$C_0\sum_j s_{j-1}b_j\le(C_0/2)S^2$.
The mixed loss terms satisfy
\[
\sum_j(H_{j-1}b_j+s_{j-1}h'_j)\le SH,
\qquad
\sum_j s_{j-1}^2b_j\le\int_0^S t^2\,dt=\frac{S^3}{3}.
\]
These are exactly \eqref{eq:sharpaggregate}.

Fix $A=314$. Since $A\ge C_0/2$, $A\ge5C_0/6$, and $A\ge5$,
we may use the simpler bounds
\begin{align}
\widehat h_v&\le\sum_w\lambda_{vw}^3\widehat h_w+AX_v^2,
\label{eq:hrec}\\
0\le D_v&\le AX_v\left(\sum_w\lambda_{vw}^3\widehat h_w+X_v^2\right).
\label{eq:drec}
\end{align}
Crucially, the propagated third moments have their exact channel factors;
the constant $A$ multiplies only newly inserted errors.

For a descendant $u$ of $v$, put
\[
\Lambda(v,u)=\prod_{e\in[v,u]}\lambda_e,
\qquad\Lambda(v,v)=1.
\]
Unfolding \eqref{eq:hrec} in \eqref{eq:drec}, and using
\eqref{eq:terminal}, gives
\begin{align}
D_v\le{}& A^2X_v
\sum_{\substack{u\succeq v\\|u|<n}}\Lambda(v,u)^3X_u^2
\notag\\
&+\frac A3X_v\sum_{\substack{u\succeq v\\|u|=n}}
\Lambda(v,u)^3X_u.
\label{eq:unfold}
\end{align}
There is no repeated multiplication by $A$ along the depth of the tree.
Each error is inserted once and thereafter propagated only by channel
products.

Multiplying \eqref{eq:deficit} by $\rho(v)$ and summing over internal
vertices telescopes, because $\rho(w)=\rho(v)\lambda_{vw}^2$:
\begin{equation}\label{eq:telescope}
\widehat m_o=M-\sum_{|v|<n}\rho(v)D_v.
\end{equation}

\subsection{The heterogeneous path identity}

Choose $V_n\in L_n$ with probability $\Theta(V_n)/M$, and let
$V_0=o,V_1,\ldots,V_n$ be its ancestors. Conservation gives
\begin{equation}\label{eq:pathlaw}
\Pp(V_j=u)=\frac{\Theta(u)}M.
\end{equation}
Every path with positive probability extends to an infinite ray and
satisfies \eqref{eq:energy}. The exact identity
\begin{equation}\label{eq:keyidentity}
\boxed{\quad \rho(v)\Lambda(v,u)^3X_u
=\Lambda(v,u)\Theta(u)\quad}
\end{equation}
is the required transfer of the third-moment propagation to the flow.
Indeed, $\rho(u)=\rho(v)\Lambda(v,u)^2$.

For fixed depths $i\le j<n$, summing over descendants in the first
term of \eqref{eq:unfold} becomes
\begin{align*}
&\sum_{|v|=i}\rho(v)X_v
\sum_{\substack{u\succeq v\\|u|=j}}\Lambda(v,u)^3X_u^2\\
&\hspace{20mm}=M\E\bigl[\Lambda(V_i,V_j)X_{V_i}X_{V_j}\bigr].
\end{align*}
The terminal term similarly becomes
$M\E[\Lambda(V_i,V_n)X_{V_i}]$.
Hence
\begin{align}
\frac1M\sum_{|v|<n}\rho(v)D_v
\le{}& A^2\E\sum_{0\le i\le j<n}
\Lambda(V_i,V_j)X_{V_i}X_{V_j}\notag\\
&+\frac A3\E\sum_{i<n}\Lambda(V_i,V_n)X_{V_i}.
\label{eq:pathloss}
\end{align}
The kernel is path-dependent, but
$\Lambda(V_i,V_j)\le q^{j-i}$. For each fixed $k\ge0$,
Cauchy--Schwarz and \eqref{eq:energy} imply pathwise
\[
\sum_i X_{V_i}X_{V_{i+k}}\le2\varepsilon^2.
\]
Also $X_{V_i}\le\varepsilon$. Therefore \eqref{eq:pathloss} is at
most
\begin{equation}\label{eq:lossbound}
\frac{2A^2\varepsilon^2}{1-q}
+\frac{Aq\varepsilon}{3(1-q)}.
\end{equation}
Choose once and for all
\begin{equation}\label{eq:eps}
\varepsilon=\frac{1-q}{2A}=\frac{1-q}{628}.
\end{equation}
Then \eqref{eq:lossbound} is at most
$(1-q)/2+q/6\le1/2$. In particular, \eqref{eq:telescope} gives
$\widehat m_o\ge M/2$ for every $n$. The original level observation
has at least this information, so
\begin{equation}\label{eq:advantage}
a_n\ge\frac{\varepsilon\alpha}{16}>0\qquad(n\ge1).
\end{equation}
This proves reconstruction.

For the quantitative lower bound, write $C=\Cap_3(T,\rho)$.
If $C>0$, an admissible flow can be chosen and scaled to mass
$\alpha=\min\{1,C/2\}$, which is at least
$\tfrac12\min\{C,1\}$. (The capacity is finite because
the root has finitely many children.) Since the original observation has
quadratic information at least $\widehat m_o\ge\varepsilon\alpha/2$,
we obtain
\[
 m_\infty\ge\frac{1-q}{2512}\min\{C,1\},
 \qquad
 a_\infty\ge\frac{1-q}{20096}\min\{C,1\}.
\]
If $C=0$, the lower bounds are immediate. This completes
Theorem~\ref{thm:main}.

\begin{proof}[Proof of Corollary~\ref{cor:robust}]
The reverse implication follows from data processing. Assume $C>0$ and use
the same admissible flow. Put
\[
 \varepsilon_\eta
 =\min\left\{\frac{1-q}{628},\,3\eta^2\right\}.
\]
Then $X_v\le\varepsilon_\eta\le3\eta^2$ at every terminal. Starting
from the already-noisy terminal spin, whose channel parameter is $\eta$,
apply the further channel of parameter
$\sqrt{X_v/3}/\eta\le1$. The resulting terminal experiment is exactly
the one used in \eqref{eq:terminal}. All subsequent estimates are unchanged,
and the loss bound is no larger than $1/2$ because
$\varepsilon_\eta\le(1-q)/628$. Thus
\[
 a_n^{(\eta)}\ge\frac{\varepsilon_\eta\alpha}{16}
 \ge\frac{\varepsilon_\eta}{32}\min\{C,1\}
\]
for every $n$, proving \eqref{eq:robust-lower} and the corollary.
\end{proof}
